% =====================================================================
%  قالب باب جديد — انسخ هذا الملف باسم جديد ثم أضفه في main.tex بـ \include
%  مثال:  chapters/05-idgham-kabir.tex   ثم   \include{chapters/05-idgham-kabir}
% =====================================================================
\chapter{اسم الباب}

مقدمة قصيرة: ما موضوع الباب؟ وكم عدد أبياته؟

\section{عنوان المسألة}

% ---------- البيت ----------
% \bayt{رقم}{الشطر الأول}{الشطر الثاني}
% داخل البيت:
%   \qw{كلمة}          ← الكلمة القرآنية المختلف فيها (تظليل أصفر)
%   \rz{Asim}{نَاصِرٌ}  ← كلمة فيها رمز قارئ (لون القارئ). المفاتيح:
%                         Nafi, IbnKathir, AbuAmr, IbnAmir, Asim, Hamza, Kisai
%   \rzg{صُحْبَةٌ}      ← رمز جماعة (ث خ ذ ظ غ ش، صحبة، صحاب، عم، سما، حق، نفر، حرمي، حصن)
%   \qname{Hamza}{حَمْزَةٌ} ← اسم صريح ملوّن داخل البيت
%   تنبيه: لوّن الكلمة كاملة، لا حرفًا واحدًا منها، حتى لا ينقطع اتصال الحروف العربية.
\bayt{000}{الشطر الأول}{الشطر الثاني}

\begin{mufradat}
\kalima{الكلمة}{معناها.}
\end{mufradat}

\maana الشرح المختصر للبيت.

\begin{qiraat}
\qira{\Asim{} و\Kisai}{القراءة الأولى: \ayah{...}}
\qira{\Baqun}{القراءة الثانية}
\end{qiraat}

\begin{qaida}[عنوان القاعدة]
تفسير القاعدة.
\end{qaida}

\begin{amthila}
\begin{itemize}
\item \ayah{نص الآية} \src{السورة: رقم الآية} \ila{} القراءة.
\end{itemize}
\end{amthila}

\begin{tahrir}
\ithaf{رقم}{شطر}{شطر}
الشرح.
\end{tahrir}

\begin{tanbih}
تنبيه.
\end{tanbih}

\section{خلاصة الباب}
\begin{khulasa}[خلاصة]
جدول أو نقاط.
\end{khulasa}

% أسماء القراء الجاهزة للنص:
% \Nafi \Qalun \Warsh | \IbnKathir \Bazzi \Qunbul | \AbuAmr \AbiAmr \DuriAbuAmr \Susi
% \IbnAmir \Hisham \IbnDhakwan | \Asim \Shuba \Hafs | \Hamza \Khalaf \Khallad
% \Kisai \AbulHarith \DuriKisai | \Baqun
